\section{序論}\label{sec:intro}

\subsection{refined large $N$ 双対と KN の正値性予想}
$U(N)$ Chern--Simons 理論の large $N$ 極限は resolved conifold $X=\mathcal{O}(-1)\oplus\mathcal{O}(-1)\to\mathbb{P}^1$ 上の位相的弦理論と等価であり，
結び目 $K$ を加えると，彩色 HOMFLY-PT 多項式は $X$ の中の Lagrangian $L_K$ に端をもつ BPS 状態の数（LMOV 不変量）で書ける
\cite{Ooguri99,Labastida01}．
Kameyama--Nawata~\cite{Kameyama17} は，この描像を refined Chern--Simons 理論~\cite{Aganagic11} に広げた．
M 理論の refined index と Schwinger の計算から自由エネルギーの形を定め，refined CS 不変量 $\rCS_\lambda(K;a,q,t)$ から
refined reformulated invariants $f^q_\mu$, $f^{\bar t}_\mu$ と $\hf_\rho$ を作る式を与えた．
そのうえで，$\hf_\rho$ を太字の変数 $(\ba,\bq,\bt)$ で書いたときの係数 $\hbN_{\rho;i,j,k}$（$G$ 展開の前）と，
$G=(\bq\bt-\bq^{-1}\bt^{-1})(\bq-\bq^{-1})$ の冪に展開した後の refined BPS 不変量 $\hN_{\rho,g,\beta,J_r}$ が，非負整数になることを予想した．
物理的には，この正値性は追加の $U(1)_R$ 対称性が BPS 状態に自明に作用することから期待される．

KN がこの予想を確かめたのは，framing 0 の次のデータである（以下「KN のデータ」と呼ぶ）．
トーラス結び目・トーラス絡み目の $|\lambda|\le 3$ の彩色，
$\mathbf{4_1},\mathbf{5_2},\mathbf{6_1},\mathbf{6_3}$ などの thin な結び目の $\yng(1),\yng(2),\yng(1,1)$（彩色 HOMFLY-PT ホモロジーの Poincaré 多項式を入力とする），
および KN が構成した $\mathbf{4_1}$ の $\yng(2,1)$ の候補である．
本論文では，KN のデータの再現は規約を定めるためだけに用い，成果としては扱わない．

\subsection{本論文の内容}
本論文の目的は，KN の枠組みを framing をもつ結び目に拡張し，KN のデータを二つの方向に超えて正整数性を検証することである．
一つは結び目の方向で，トーラス結び目と thin な結び目から，代数的な iterated torus 結び目，結び目表の全ての結び目，ツイスト結び目・pretzel 結び目の族，thick な結び目へ広げる．
もう一つは framing の方向で，$x=0$ から，有限の範囲の $x$，さらに全ての整数 $x$ へ広げる．

\paragraph{framed knot への拡張（\S\ref{sec:framing}）．}
framing $x$ の不変量を $\rCS^{(x)}_\lambda=\fr[\lambda]^x\,\rCS_\lambda$，
\[
\fr[\lambda]=\bigl(-(q/t)^{1/2}\bigr)^{|\lambda|}\,q^{n(\lambda')}\,t^{-n(\lambda)}
\]
で定める．この因子は Taki~\cite{Taki07} の refined topological vertex の framing の因子の逆数にちょうど一致し，
$t=q$ では Mariño--Vafa~\cite{Marino01} の framing の規則 $W_R\mapsto(-1)^{\ell p}q^{p\kappa_R/2}W_R$ に $x=p$ で戻る．
太字の変数では $\fr[\lambda]$ は符号のない単項式 $\bt^{|\lambda|+2n(\lambda')}\bq^{2(n(\lambda')-n(\lambda))}$ になり，
framing はホモロジーの次数のずらしとして働く．
主な構造的な結果は次のとおりである．
\begin{itemize}
\item $\bt^{-|\rho|x}\hf^{(x)}_\rho$ は $x$ の指数和である（命題~\ref{prop:expsum}）．
\item 転置対称性，Laurent 性，$G$ 展開が閉じる条件は $x$ によらない（命題~\ref{prop:x-indep}）．
\item Mariño--Vafa の符号 $(-1)^{|\lambda|x}$ を外すと，奇数の $x$ で 2 箱の $\hf$ が Laurent でなくなる（命題~\ref{prop:sign}）．
\item KN の単項式の正規化のうち彩色によって変わる部分は，ちょうど Seifert framing の分の framing である（命題~\ref{prop:N2-framing}）．
\end{itemize}

\paragraph{検証（\S\ref{sec:one}--\S\ref{sec:three}）．}
入力の確からしさは結び目のクラスによって異なる（表~\ref{tab:overview}）．
\begin{itemize}
\item 基本表現 $\rho=\yng(1)$：入力が非負係数で DGR の対称性をもてば，$G$ 展開は整数で閉じ，前後とも非負である（命題~\ref{prop:one}）．
DGR の対称性は Gorsky--Hogancamp--Mellit~\cite{Gorsky21} の定理なので，ホモロジーを入力とする限り，任意の結び目・任意の framing で $\yng(1)$ は正である．
\item 2 箱：$\hf_{\yng(2)}$ が Laurent であるためには $\dim H^{\yng(2)}=(\dim H^{\yng(1)})^2$ が必要である（命題~\ref{prop:two-necessary}）．
GNSSS~\cite{Gukov15} の予想式を入力として，ツイスト結び目の全ての撚りの数で 2 箱が正であることを示す（定理~\ref{thm:twist}）．
また，8 個の結び目では全ての整数の framing で 2 箱が正である（定理~\ref{thm:two-all-x}）．
\item 3 箱：トーラス結び目 7 個と，代数的な iterated torus 結び目 7 個で，$|\rho|\le3$ の全ての $\rho$ について $G$ 展開の前後とも正である（計算結果~\ref{cr:three}）．
後者は Cherednik--Danilenko~\cite{Cherednik14} の DAHA による定義どおりの refined 不変量を入力とし，
3 本ケーブル $K7$ と 2 段の反復 $K10$（3 箱で約 183 万項）を含む．
入力の $\rCS_{\yng(2,1)}$ 自体は負の係数をもつが，BPS の段では正になる．
\end{itemize}
これらに基づき，framed knot の refined LMOV 不変量の正値性を予想~\ref{conj:A} として，
代数的な iterated torus 結び目の正値性を予想~\ref{conj:B} として提出する．
双曲結び目の 3 箱には，定義に基づく $\yng(2,1)$ の refined 不変量が無い．構造的な条件では $\yng(2,1)$ が一意に決まらないことを \S\ref{sec:21} で述べる．

\begin{table}[t]
\centering
\caption{検証の全体像．○ は $G$ 展開の前後とも正，— は未実施，× は入力が無い．}\label{tab:overview}
\resizebox{\textwidth}{!}{%
\begin{tabular}{llccccc}
\toprule
結び目 & 入力 & $\yng(1)$ & 2 箱 $x=0$ & 2 箱 $x\neq0$ & 3 箱 $x=0$ & 3 箱 $x\neq0$ \\
\midrule
自明な結び目 & -- & 全 $x$ & ○ & ○ & ○ & $|x|\le6$ \\
$T_{2,3},T_{2,5},T_{2,7},T_{3,4}$ & EHA & 全 $x$ & ○ & ○ & ○ & $|x|\le 3$--$6$ \\
$T_{3,5},T_{2,11},T_{3,7},T_{4,5}$ & EHA & 全 $x$ & ○ & $|x|\le10$ & ○ & $|x|\le10$ \\
$K1,K3$ & DAHA & 全 $x$ & ○ & ○ & ○ & $|x|\le2$（$K1$ は $x=\pm6$ も） \\
$K6a,K7,K8,K9$ & DAHA & 全 $x$ & ○ & $|x|\le10$ & ○ & $|x|\le10$ \\
$K10$ & DAHA & 全 $x$ & ○ & $|x|\le10$ & ○ & $|x|\le2$ \\
$\mathbf{4_1}$ & ホモロジー（$\yng(2,1)$ は KN の候補） & 全 $x$ & ○ & ○ & ○ & $|x|\le4$ \\
結び目表 801 個 & ホモロジー & 全 $x$ & × & × & × & × \\
ツイスト $K_p$ & GNSSS の予想式 & 全 $x$ & 全 $p$ & 注 1 & × & × \\
pretzel $K_{s,t}$ & GNSSS の予想式 & 全 $x$ & 注 2 & -- & × & × \\
$\mathbf{10_{139}},\mathbf{10_{152}}$ & quiver ＋ KRSS の予想式 & 全 $x$ & ○ & -- & × & × \\
\bottomrule
\end{tabular}}

\smallskip
\raggedright
注 1：$p$ と $x$ が同符号の向きでは全て正．逆符号は $|p|\le6$，$|x|\le14$ の直接計算で正．\\
注 2：同符号は全て正．異符号は $\min(|s|,t)\le20$ で証明，$|s|\le30$ かつ $t\le40$ で確認．
\end{table}

\subsection{関連する研究}
非 refined の framed knot の LMOV 型の整数性は Mariño--Vafa~\cite{Marino01} に始まり，
Zhu らの一連の研究~\cite{Zhu21,Luo16,Luo21,Wang25,Zhu25} がある．
Zhu~\cite{Zhu21} の ``refined LMOV conjecture'' の refinement は KN の refined（$t\neq q$）とは別の意味であり（同論文の Remark），本論文の主張とは対象が異なる．
refined の open BPS 不変量の非負性は，単一のブレーン・対称表現について Cheng--Sułkowski~\cite{Cheng21}（framing は固定）と
Banerjee--Ishtiaque--Jeong~\cite{Banerjee26} が論じている．
knots--quivers 対応~\cite{Kucharski17,Kucharski17b,Ekholm18} では，framing の変化は quiver 行列のずらし $C\mapsto C+f$ として現れ，
対称表現の refined LMOV 不変量の整数性が quiver の構造から従う．
本論文の調べた範囲では，KN の $\hf_\rho$ を framing をつけて作り，非対称な $\rho$（$\yng(2,1)$ を含む）まで $G$ 展開の前後の正整数性を調べた研究は見当たらない．
ケーブルの refined 不変量の定義は Cherednik--Danilenko~\cite{Cherednik14,Cherednik15} による．
その明示例は $K1$ の $\yng(1),\yng(2)$ などに限られ，$|\lambda|=3$ の値と refined BPS 不変量は文献に見当たらない．

\subsection{構成}
\S\ref{sec:prelim} で変数と KN の refined LMOV 不変量を復習する．
\S\ref{sec:framing} で framed knot への拡張と，$x$ によらない構造を述べる．
\S\ref{sec:inputs} で入力（トーラス結び目，ケーブル，ホモロジー，予想式）と計算法を説明する．
\S\ref{sec:one}，\S\ref{sec:two}，\S\ref{sec:three} で，それぞれ $\yng(1)$，2 箱，3 箱の検証を述べる．
\S\ref{sec:discussion} で未解決の問題を議論する．
付録に計算法の詳細，非代数的なケーブル，文献の誤植，データとコードをまとめる．

主張の確からしさを区別するため，解析的に導いたものを定理・命題・補題とし，計算機で確かめたものは「計算結果」として範囲とともに書く．
予想式を入力とする主張は，その旨を文中に明記する．
